All evaluation is performed on held-out test cases only. For each task the model
is trained on a single source contrast and evaluated on every contrast of that
task's test split, including the training contrast itself; training and
validation subjects never appear in any reported number.

\paragraph{ON-Harmony} \cite{warrington2025onharmony} --- 31-region healthy-brain
segmentation. Reference labels are \emph{silver standard}: they were produced
automatically with FreeSurfer~\cite{fischl2012freesurfer} on the T1w volume of
each subject and then registered to that subject's other contrasts, which are
acquired in the same session. They are therefore not manual annotations, and
absolute Dice on this task should be read as agreement with an automated
reference rather than with a human rater; the comparison between methods, which
all use the identical reference, is unaffected.

\paragraph{Open-MS} \cite{lesjak2018openms} --- MS-lesion segmentation with
manual consensus annotations, evaluated across FLAIR, T1w and T2w.

\paragraph{Brats2024-GLI} \cite{deverdier2024brats2024} --- post-treatment
glioma sub-region segmentation, evaluated across T1n, T1c, T2w and T2-FLAIR.

\paragraph{CHAOS} \cite{kavur2021chaos} --- abdominal organ segmentation (liver,
spleen, both kidneys), evaluated across the CHAOS MR contrasts and, as a
cross-modality transfer, on held-out CT from AMOS~\cite{ji2022amos} and
SLIVER07~\cite{heimann2009comparison}.

\paragraph{ToothFairy2} \cite{bolelli2025toothfairy2} --- mandible segmentation
in dental cone-beam CT, the one non-CT, non-MRI training modality here. The
released label set is 42 classes, but annotation coverage is cohort-dependent:
several structures are annotated in only part of the cohort, so scoring them
would measure which cohort a case came from as much as the model. We therefore
reduce to the mandibular structures annotated throughout, and every column of
this task --- in-domain and cross-dataset alike --- is scored on that same
common label, so the numbers being averaged are commensurable. ToothFairy2
acquires CBCT only, so its held-out evidence is entirely
cross-modality: HaN-Seg~\cite{podobnik2023hanseg} head-and-neck CT and MR-T1.
Only HaN-Seg's CT arm is usable as a paired test; its MR is not registered to
its CT, and the MR-T1 column is scored against its own MR annotation.

\paragraph{I-SPY2 and Duke-Breast-MRI} \cite{li2022ispy2,saha2018duke} ---
breast-lesion segmentation. I-SPY2 (CC BY 4.0) is a neoadjuvant-chemotherapy
trial whose enrolment requires biopsy-proven invasive cancer, and its reference
is the trial's own functional tumour volume mask; training uses its
first-post-contrast T1 (T1-CE) and its T2w. Duke-Breast-MRI, obtained through
the MAMA-MIA release~\cite{garrucho2025mamamia} (CC BY-NC 4.0, so it carries a
non-commercial term the other cohorts here do not), is an external
cross-dataset test cohort only, contributing a T1-CE and a pre-contrast column.
Both are malignant-only cohorts with expert tumour masks. I-SPY2's DCE series
are acquired unilaterally at some sites and bilaterally at others; both
field-of-view variants are generated per patient rather than discarding either,
which is why the two training contrasts yield different test-case counts (102
for T1-CE, 168 for T2w) from the same 84 held-out patients.

\begin{table}[h]
\centering
\small
\setlength{\tabcolsep}{4pt}
\begin{tabular}{lcc}
\toprule
Task & Train & Test \\
\midrule
CHAOS          & 16  & 24 \\
ON-Harmony     & 84  & 8  \\
Open-MS        & 22  & 8  \\
Brats2024-GLI  & 630 & 70 \\
ToothFairy2    & 408 & 71 \\
I-SPY2         & 476 & 84 \\
\bottomrule
\end{tabular}
\caption{Per-task subject/case counts, pooled across the 3 folds used
everywhere in this paper ($\pm 1$ between folds). ON-Harmony counts are scan
sessions, not unique subjects, since its travelling-heads design scans some
subjects at multiple sites. CHAOS test includes its own held-out CT split,
separate from the AMOS/SLIVER07 cross-dataset transfer counted
independently (external cohorts, not re-tabulated here); likewise
ToothFairy2's HaN-Seg transfer and I-SPY2's Duke-Breast-MRI transfer. I-SPY2
counts are patients: each held-out patient yields one T1-CE case and, where
both field-of-view variants exist, more than one T2w case, giving the 102 and
168 per-contrast case counts quoted above. ToothFairy2 excludes one case with
no mandible label.}
\label{tab:suppl-datasets}
\end{table}
